\chapter{Market-Data Distribution and Entitlements}\label{ch:pl:market-data-distribution-and-entitlements}

An exchange's auditor asks a firm for the list of every person and every program that used the exchange's depth data in the past three years.
The firm can name forty screens: it pays for forty. It cannot name, at first, the eleven applications that read the same feed from the
firm's distribution platform --- strategies, a risk service, a tick recorder, a surveillance system, a research dashboard --- and each of them is
a use the exchange charges for. The firm in this chapter is a model, but the documents it is audited against are real: one exchange's price
list, data policies and subscriber agreement. The chapter builds the platform that distributes market data inside a firm and the part of it
that is usually missing: the record of who used what, for which purpose, counted the way the venue counts.

\section{Distributing data inside the firm}

A direct feed (Book~1, chapter~9) arrives once at the firm's feed handlers. From there it must reach dozens of consumers with very different
needs: a strategy that wants every message within microseconds, a risk service that wants current prices, a trader's screen that cannot show
more than a few updates a second, a research job that records everything. Connecting each of them to the feed would multiply the network
load, the licences and the points of failure; the firm puts one platform in between.

\begin{definition}[Market-data distribution platform]\label{def:pl:market-data-distribution-and-entitlements:platform}
A \emph{market-data distribution platform}\index{market-data distribution platform} receives normalised market data once and delivers it
inside the firm to subscribers by topic (an instrument, a venue, a data type), with a delivery policy per subscriber --- every message in
order, or conflated --- and with an \omterm{def:pl:market-data-distribution-and-entitlements:entitlement}{entitlement} check on every subscription.
\end{definition}

\begin{omfigure}
\begin{tikzpicture}[font=\small,node distance=5mm,
  box/.style={draw=omDef,fill=omDef!8,rounded corners=2pt,align=center,minimum height=9mm,text width=22mm},
  chk/.style={draw=omThm,fill=omThm!8,rounded corners=2pt,align=center,minimum height=9mm,text width=22mm},
  arr/.style={-{Stealth[length=2mm]},thick}]
\node[box] (f) {feed handler\\(normalised)};
\node[box,right=8mm of f,minimum height=34mm,text width=18mm] (b) {bus:\\topics,\\one queue per\\subscriber};
\node[box,right=12mm of b,yshift=19.5mm,text width=36mm,minimum height=7mm] (s1) {strategy, every message};
\node[box,right=12mm of b,yshift=6.5mm,text width=36mm,minimum height=7mm] (s2) {risk service};
\node[box,right=12mm of b,yshift=-6.5mm,text width=36mm,minimum height=7mm] (s3) {research dashboard};
\node[box,right=12mm of b,yshift=-19.5mm,text width=36mm,minimum height=7mm] (s4) {trader's screen (conflated)};
\node[chk,below=12mm of f] (e) {entitlement\\system};
\node[chk,right=4mm of e,xshift=22mm,yshift=-12mm] (u) {usage log};
\node[chk,right=6mm of u] (r) {usage report\\and audit};
\draw[arr] (f) -- (b);
\foreach \x in {s1,s2,s3,s4} \draw[arr] (b.east) -- (\x.west);
\draw[arr,omThm] (e.east) -- node[above,font=\scriptsize,sloped]{check} (b.south west);
\draw[arr,omThm] (b.south) |- (u.west);
\draw[arr,omThm] (u) -- (r);
\end{tikzpicture}

\omcaption{The chapter's distribution platform: one normalised feed, one bus with a queue per subscriber, an \omterm{def:pl:market-data-distribution-and-entitlements:entitlement}{entitlement} check at subscription
and at delivery, and every access written to a usage log from which the monthly \omterm{def:pl:market-data-distribution-and-entitlements:report}{usage report} and the audit \omterm{def:pl:post-trade-systems:recon}{reconciliation} are
produced.}\label{fig:pl:market-data-distribution-and-entitlements:platform}
\end{omfigure}

The chapter's bus (\texttt{firm.entitle}) is simulated in virtual time. Ten minutes of depth updates for twenty instruments --- 240\,267
messages, from Book~7's tape generator, the order flow that Book~10's simulator replays --- go to four subscribers: a strategy that takes
20~microseconds per message, a risk service that takes 2~milliseconds, a research dashboard that takes 10, and a trader's screen following
five instruments at 50~milliseconds a refresh. Each subscriber has its own queue, so that a slow one cannot delay the others: the
publish--subscribe design of chapter~\ref{ch:pl:messaging-and-event-driven-architecture}, and the slow-consumer problem of Book~13,
chapter~12, one level up.

\omcode{code/firm/entitle/firm_entitle.py}{98}{111}{The two delivery policies: a queue of every update in order, or conflation, which keeps only
the latest update per topic and serves the topic that has waited longest.}\label{lst:pl:market-data-distribution-and-entitlements:queue}

\begin{table}[htbp]
\centering\small
\begin{tabular}{llrrrr}
\toprule
policy & subscriber & delivered & largest queue & mean age & largest age\\
\midrule
every message & strategy & 240\,267 & 1\,199 & 0.1~ms & 24~ms\\
 & risk service & 240\,267 & 1\,203 & 58~ms & 2.4~s\\
 & research dashboard & 240\,267 & 180\,267 & 891~s & 1\,803~s\\
conflated & strategy & 238\,870 & 20 & 0.0~ms & 0.4~ms\\
 & risk service & 221\,364 & 20 & 4.0~ms & 32~ms\\
 & research dashboard & 60\,012 & 20 & 67~ms & 134~ms\\
\bottomrule
\end{tabular}
\caption{Ten minutes of depth updates for twenty instruments delivered to three subscribers, with every message queued or conflated per
instrument. The age of a message is the time from its publication to the end of its processing. The screen, conflated in both runs, sees a
mean age of 125~ms. The strategy's largest queue is the opening snapshot, which arrives at one instant.}
\label{tab:pl:market-data-distribution-and-entitlements:policies}
\end{table}

The research dashboard is the lesson of \cref{tab:pl:market-data-distribution-and-entitlements:policies}. It can process 100 messages a
second and the feed brings about 400: with every message queued, its queue grows through the whole session to 180\,267 messages, it shows prices
fifteen minutes old on average, and it is still working through the session half an hour after it ended
(\cref{fig:pl:market-data-distribution-and-entitlements:age}). Conflated, it processes a quarter of the messages and never shows a price
older than 134~ms. The risk service can keep up on average but not in bursts; conflation takes its worst age from 2.4~seconds to 32~ms. A
screen or a risk service wants the current state, not the history, and should be conflated; a strategy or a recorder needs every message
and must be fast enough not to need it.

\begin{omfigure}
\begin{tikzpicture}
\begin{semilogyaxis}[width=11.5cm,height=5.6cm,xlabel={time since the start of the session (s)},ylabel={age of the price shown (s, log scale)},
  grid=major,grid style={gray!20},tick label style={font=\small},label style={font=\small},table/col sep=comma,xmin=0,xmax=2450,
  ymin=0.01,ymax=5000,xtick={0,600,1200,1800,2400}]
\addplot[omThm,thick,mark=none] table[x=t,y=age_s]{figdata/platforms/23-market-data-distribution-and-entitlements/age_all.csv};
\addplot[omDef,thick,mark=none] table[x=t,y=age_s]{figdata/platforms/23-market-data-distribution-and-entitlements/age_conflate.csv};
\draw[gray,dashed] (axis cs:600,0.01) -- (axis cs:600,5000) node[pos=0.5,right,font=\scriptsize]{session ends};
\node[omThm,font=\scriptsize,anchor=north] at (axis cs:1500,250) {every message queued};
\node[omDef,font=\scriptsize,anchor=south west] at (axis cs:20,0.15) {conflated per instrument};
\end{semilogyaxis}
\end{tikzpicture}

\omcaption{Age of the prices shown by the research dashboard, sampled every 2\,000 deliveries (queued) or 500 (conflated). Queued, the age
grows through the session and the backlog takes another half hour to drain; conflated, it stays near a tenth of a second. Data:
\texttt{fig\_entitle.py}.}\label{fig:pl:market-data-distribution-and-entitlements:age}
\end{omfigure}

\section{Licences, uses and units of count}

What the firm pays for is not the feed: it is each use of it. Venues distinguish display use, by a person looking at a screen, from
non-display use by a machine, which Book~1, chapter~29, introduced with the non-display fee; and they count each in a stated unit.

\begin{definition}[Unit of count]\label{def:pl:market-data-distribution-and-entitlements:unit}
A \emph{unit of count}\index{unit of count} is what a venue's licence counts to charge for its data: for display use typically each device or
each unique user identification, so that one person with two logins counts twice; for non-display use typically each device (usually a
server) that receives and benefits from the data, or each person who can change the application in real time.
\end{definition}

The definition is the venue's, and it is where firms lose money. The chapter's venue --- Nasdaq, whose documents are in the ledger --- counts a
display subscriber as \textquotedbl{}a device, computer terminal, automated service, or unique user identification and password
combination\textquotedbl{}, and non-display use by the servers that receive and benefit from the data, including those that compute on it or
create \omterm{def:pl:market-data-distribution-and-entitlements:derived}{derived data} from it. Servers that only transport, normalise or store the data are not necessarily fee-liable, but the firm must be able
to show which they are. A firm that counts people instead of identifiers, or applications instead of servers, undercounts by construction.

\section{Entitlements}

\begin{definition}[Entitlement, entitlement system]\label{def:pl:market-data-distribution-and-entitlements:entitlement}
An \emph{entitlement}\index{entitlement} is a permission for a subject --- a user identifier or an application --- to receive one source of
data for one use. An \emph{entitlement system}\index{entitlement system} holds the entitlements with their grant and revocation times,
answers whether a subject is entitled at a given time, and is checked by the distribution platform before any data is delivered.
\end{definition}

The question is when to check. Checking at subscription is cheap and enough while nothing changes; but a trader who leaves the desk, or an
application whose licence lapses, keeps receiving until it next subscribes. In the chapter's run the screen user's \omterm{def:pl:market-data-distribution-and-entitlements:entitlement}{entitlement} is revoked five
minutes into the session. With the check made only at subscription, the screen receives another 6\,005 updates it was not entitled to; checked
at every delivery, none.

\omcode{code/firm/entitle/firm_entitle.py}{133}{147}{Delivery: the entitlement is checked again at each delivery, so that a revocation takes
effect at once, and the age of each delivered message is recorded.}\label{lst:pl:market-data-distribution-and-entitlements:serve}

A real platform does not call the \omterm{def:pl:market-data-distribution-and-entitlements:entitlement}{entitlement system} for each of millions of messages: it caches the answer per subscription and invalidates the
cache when an \omterm{def:pl:market-data-distribution-and-entitlements:entitlement}{entitlement} changes, which has the same effect at almost no cost. What matters is that revocation is pushed to the platform
rather than discovered at the next login.

\section{Usage reporting and exchange audits}

\begin{definition}[Usage report]\label{def:pl:market-data-distribution-and-entitlements:report}
A \emph{usage report}\index{usage report} is the periodic declaration a firm makes to a data source of how many units of count used its data,
by product and by use, from which the source bills its fees.
\end{definition}

\begin{definition}[Market-data audit]\label{def:pl:market-data-distribution-and-entitlements:audit}
A \emph{market-data audit}\index{market-data audit} is a venue's or vendor's review of a firm's actual use of its data over a past period,
reconciled with what the firm reported; the difference is billed back, usually with interest and, above a threshold, the audit's costs.
\end{definition}

The chapter's firm reported 40 display subscribers every month for three years and no non-display use. Its usage log --- which the platform
writes on every subscription --- tells a different story (\cref{fig:pl:market-data-distribution-and-entitlements:usage}). Traders grew from 38
to 47, and six of them have a second login, so display identifiers went from 43 to 53; the eleven applications came into use one after another
and by the end ran on 22 servers.

\omcode{code/firm/entitle/firm_entitle.py}{168}{195}{The usage report counts distinct units of count by use; the audit reconciles each month
with the licences held and prices the difference with the fee schedule.}\label{lst:pl:market-data-distribution-and-entitlements:audit}

\begin{omfigure}
\begin{tikzpicture}
\begin{axis}[width=11.5cm,height=5.6cm,xlabel={month},ylabel={units of count},
  grid=major,grid style={gray!20},tick label style={font=\small},label style={font=\small},table/col sep=comma,
  xmin=0,xmax=35,ymin=0,ymax=63,const plot]
\addplot[omDef,thick] table[x=month,y=display_used]{figdata/platforms/23-market-data-distribution-and-entitlements/usage.csv};
\addplot[omDef,thick,dashed] table[x=month,y=display_licensed]{figdata/platforms/23-market-data-distribution-and-entitlements/usage.csv};
\addplot[omThm,thick] table[x=month,y=nd_used]{figdata/platforms/23-market-data-distribution-and-entitlements/usage.csv};
\node[omDef,font=\scriptsize,anchor=south] at (axis cs:20,53) {display identifiers used};
\node[font=\scriptsize,anchor=north] at (axis cs:24,39) {display subscribers reported: 40};
\node[omThm,font=\scriptsize,anchor=south] at (axis cs:25,24) {non-display servers used (reported: 0)};
\end{axis}
\end{tikzpicture}

\omcaption{Three years of the model firm's usage log against what it reported. Display identifiers exceed the 40 reported from the first month
(43, because of second logins) and reach 53; non-display servers grow from 4 to 22 as the eleven applications come into use. Data:
\texttt{fig\_entitle.py}.}\label{fig:pl:market-data-distribution-and-entitlements:usage}
\end{omfigure}

Priced with the venue's 2026 fees for its depth product --- \$84 a month per professional display subscriber, and for non-display use \$412 a
month per subscriber up to 39, then a flat \$16\,490 a firm from 40 to 99 (\cref{fig:pl:market-data-distribution-and-entitlements:fees}) ---
the shortfall is \$1\,900 in the first month and \$10\,156 in the last, \$258\,596 over the three years: \$22\,932 for display and \$235\,664 for
non-display. The unreported units reach 87.5\% of those reported, far above the 10\% at which the venue's subscriber agreement adds the
audit's costs; interest comes on top, at a rate the model does not price. The display shortfall is the smaller part. The firm's real mistake
was believing that data it had paid to receive could be used by any program.

\begin{omfigure}
\begin{tikzpicture}
\begin{axis}[width=11.5cm,height=5.4cm,xlabel={non-display subscribers (servers)},ylabel={monthly fee (\$ thousand)},
  grid=major,grid style={gray!20},tick label style={font=\small},label style={font=\small},table/col sep=comma,
  xmin=0,xmax=120,ymin=0,ymax=38,const plot]
\addplot[omThm,thick] table[x=devices,y expr=\thisrow{fee}/1000]{figdata/platforms/23-market-data-distribution-and-entitlements/fees.csv};
\draw[omDef,dashed] (axis cs:22,0) -- (axis cs:22,9.064);
\draw[omDef,-{Stealth[length=1.5mm]}] (axis cs:55,6) node[right,font=\scriptsize]{the model firm: 22 servers, \$9\,064} -- (axis cs:23,8.8);
\node[font=\scriptsize,anchor=south east] at (axis cs:37,17.5) {39: \$16\,068};
\node[font=\scriptsize,anchor=south] at (axis cs:70,17.5) {40 to 99: \$16\,490};
\node[font=\scriptsize,anchor=north east] at (axis cs:98,32) {100 to 249: \$32\,990};
\end{axis}
\end{tikzpicture}

\omcaption{The venue's 2026 monthly non-display fee for its depth product against the number of fee-liable subscribers: \$412 each up to 39,
then flat per firm. Between 40 and 99 servers the fee does not change, which is why the count and the tier, not the price, drive the bill.
Data: \texttt{fig\_entitle.py}.}\label{fig:pl:market-data-distribution-and-entitlements:fees}
\end{omfigure}

The remedy is technical as much as contractual. The platform's usage log is the evidence an audit asks for, and it must record the \omterm{def:pl:market-data-distribution-and-entitlements:unit}{unit of
count} the venue uses --- identifiers, not people; servers, not applications --- for as long as the look-back. The chapter's venue limits
liability for a good-faith error to three years before the underreporting was found, so three years of logs is the minimum. A monthly
\omterm{def:pl:post-trade-systems:recon}{reconciliation} of the log against the licences, run by the firm before any auditor does, turns a back-bill into a licence purchase.

\section{Derived data}

\begin{definition}[Derived data]\label{def:pl:market-data-distribution-and-entitlements:derived}
\emph{Derived data}\index{derived data} is information computed in whole or in part from a venue's data that cannot be reverse-engineered to
recreate it or used as a reasonable substitute for it: an index, a signal, a model's fair value. What is and is not derived data is decided by
each venue's policy.
\end{definition}

A conflated quote is not \omterm{def:pl:market-data-distribution-and-entitlements:derived}{derived data}: it is the venue's data, delivered less often, and needs the same licence. A mid-price published every
second may be a substitute for the quote; a volatility estimate or a basket's fair value usually is not. The distinction matters in two
directions: \omterm{def:pl:market-data-distribution-and-entitlements:derived}{derived data} may be redistributed under terms that the raw data does not allow, and the servers that create it are, for the
chapter's venue, among the non-display units that count. The \omterm{def:pl:market-data-distribution-and-entitlements:entitlement}{entitlement} model should therefore carry the use --- display, non-display,
derived --- and not only the source, and the \omterm{def:pl:a-tick-store-on-open-formats:store}{tick store} of chapter~\ref{ch:pl:a-tick-store-on-open-formats} should record which licence
each dataset was captured under, since the licence follows the data into research.

\begin{dated}{2026-09}{One venue's data policies and fees}\label{dat:pl:market-data-distribution-and-entitlements:nasdaq}
Nasdaq's US equities price list sets, from 1 January 2026, monthly fees for Nasdaq TotalView of \$84.00 per professional subscriber (\$80.50
in 2025) and \$15.00 per non-professional; for non-display use of its depth data with direct access, \$412.00 per subscriber from 1 to 39,
\$16\,490 per firm from 40 to 99, \$32\,990 from 100 to 249 and \$75\,000 from 250; and distributor fees of \$1\,690 per firm for internal
distribution and \$3\,340 for direct access. Its US Equities and Options Data Policies (version 2.6) define the subscriber, non-display use and
\omterm{def:pl:market-data-distribution-and-entitlements:derived}{derived data} as quoted in this chapter, and state that undeclared usage of a system cannot be netted over the audit period. Its subscriber
agreement limits liability for a good-faith error to three years of unpaid fees with interest, and adds the audit's costs when underreporting
reaches 10\% of the reported units.
\end{dated}

\section{Tutorial: distribute, entitle, report, audit}\label{tut:pl:market-data-distribution-and-entitlements:tut}

\begin{tutorial}
\textbf{Goal.} Distribute a feed to subscribers of different speeds, enforce \omterm{def:pl:market-data-distribution-and-entitlements:entitlement}{entitlements}, and reconcile three years of usage with the
licences held. \textbf{End state:} \cref{tab:pl:market-data-distribution-and-entitlements:policies} and the audit's
\$258\,596.
\begin{tutsteps}
\item \textbf{Feed}: \texttt{pl\_entitle.updates()}, ten minutes of twenty instruments.
\item \textbf{Distribution}: \texttt{distribute(ups, \textquotedbl all\textquotedbl)} and
\texttt{distribute(ups, \textquotedbl conflate\textquotedbl)}; compare queues and ages.
\item \textbf{\omterm{def:pl:market-data-distribution-and-entitlements:entitlement}{Entitlements}}: rerun with \texttt{check\_on\_delivery=False} and count the screen's deliveries after the revocation.
\item \textbf{Usage}: \texttt{usage\_log()} and \texttt{firm\_entitle.usage\_report} for any month.
\item \textbf{Audit}: \texttt{audit()}, the months, the shortfalls and the amount owed.
\end{tutsteps}
\textbf{What to change next.} Move the research dashboard onto the servers of a distribution tier and see what the non-display count becomes;
price the audit with each year's fees instead of 2026's.
\end{tutorial}

\section{Build: entitlements}\label{bld:pl:market-data-distribution-and-entitlements:entitle}

\begin{build}
\bfield{Purpose} Deliver market data inside the firm to each subscriber in the form it can use, only to those entitled, with the evidence of
every use counted as the venue counts it.
\bfield{Interface} \texttt{Entitlement}, \texttt{EntitlementSystem} (\texttt{grant}, \texttt{revoke}, \texttt{check}), \texttt{Bus}
(\texttt{subscribe}, \texttt{run}), \texttt{UsageLog}, \texttt{usage\_report}, \texttt{audit}.
\bfield{Rules} One queue per subscriber; conflation for subscribers that want state; \omterm{def:pl:market-data-distribution-and-entitlements:entitlement}{entitlement} checked at subscription and enforced at
delivery; every use logged in the venue's \omterm{def:pl:market-data-distribution-and-entitlements:unit}{unit of count}; logs kept at least as long as the look-back; a monthly \omterm{def:pl:post-trade-systems:recon}{reconciliation} before any
audit.
\bfield{Acceptance tests} \texttt{code/firm/entitle/tests/}: checks by subject, source, use and time; queued and conflated delivery; a busy
subscriber scheduled correctly; revocation at delivery against at subscription; usage counts by identifier and device, and an audit's
shortfall and amount.
\bfield{Stretch} Tiered fees with enterprise licences; a derived-data register; a \omterm{def:pl:market-data-distribution-and-entitlements:report}{usage report} in a venue's file format.
\end{build}

\begin{omsources}
\item Nasdaq, \emph{U.S. Equities Price List}, fees effective 2025, 2026 and 2027.
\item Nasdaq, \emph{US Equities and Options Data Policies}, version 2.6.
\item Nasdaq, \emph{Global Subscriber Agreement}, sections 5(g) and 5(h).
\item One Quant Book~1, chapters~9 and~29 (direct feeds, non-display fees); Book~13, chapter~12 (slow consumers and conflation).
\end{omsources}

\section{Exercises}

\begin{exercise}[$\star$]\label{exo:pl:market-data-distribution-and-entitlements:1}
Why does the platform keep one queue per subscriber rather than one per topic?
\end{exercise}

\begin{exercise}[$\star$]\label{exo:pl:market-data-distribution-and-entitlements:2}
A trader uses the same depth data on a desktop and, with a second login, on a laptop. How many display subscribers is that for the chapter's
venue?
\end{exercise}

\begin{exercise}[$\star$]\label{exo:pl:market-data-distribution-and-entitlements:3}
What is the model firm's non-display fee in the last month, and what would it be with 45 servers?
\end{exercise}

\begin{exercise}[$\star\star$]\label{exo:pl:market-data-distribution-and-entitlements:4}
Which of the four subscribers should be conflated, and why not the strategy?
\end{exercise}

\begin{exercise}[$\star\star$]\label{exo:pl:market-data-distribution-and-entitlements:5}
Why does the research dashboard, queued, show prices fifteen minutes old on average when it takes 10~milliseconds per message?
\end{exercise}

\begin{exercise}[$\star\star$]\label{exo:pl:market-data-distribution-and-entitlements:6}
The firm wants to check \omterm{def:pl:market-data-distribution-and-entitlements:entitlement}{entitlements} at every delivery without calling the \omterm{def:pl:market-data-distribution-and-entitlements:entitlement}{entitlement system} millions of times a minute. How?
\end{exercise}

\begin{exercise}[$\star\star\star$]\label{exo:pl:market-data-distribution-and-entitlements:7}
\emph{Coding.} Rerun \texttt{audit} with the firm reporting its display identifiers correctly every month but still no non-display use. How
much is owed?
\end{exercise}

\begin{exercise}[$\star\star\star$]\label{exo:pl:market-data-distribution-and-entitlements:8}
\emph{Find the flaw.} \textquotedbl{}We pay the distributor fee for internal distribution, so any application inside the firm may use the
data.\textquotedbl{}
\end{exercise}

\section{Problem: Eleven Applications}

\begin{problem}[{Weekend problem --- eleven applications}]\label{pb:pl:market-data-distribution-and-entitlements:1}
The chapter's platform, its four subscribers, and the model firm's three years of usage.

\medskip\noindent\textbf{Part I --- Distribution.}
\begin{enumerate}
\item What does the distribution platform do that connecting each consumer to the feed does not?
\item How fast does each subscriber process messages, and how fast does the feed arrive?
\item What happens to the research dashboard with every message queued?
\item What does conflation change for the risk service and the dashboard?
\item What is the strategy's largest queue, and why?
\end{enumerate}

\medskip\noindent\textbf{Part II --- \omterm{def:pl:market-data-distribution-and-entitlements:entitlement}{Entitlements}.}
\begin{enumerate}[resume]
\item What does an \omterm{def:pl:market-data-distribution-and-entitlements:entitlement}{entitlement} name?
\item What happens after a revocation when \omterm{def:pl:market-data-distribution-and-entitlements:entitlement}{entitlements} are checked only at subscription?
\item How does a platform enforce \omterm{def:pl:market-data-distribution-and-entitlements:entitlement}{entitlements} at every delivery cheaply?
\item What is the \omterm{def:pl:market-data-distribution-and-entitlements:unit}{unit of count} for display and for non-display use at the chapter's venue?
\item Which servers are not necessarily fee-liable?
\end{enumerate}

\medskip\noindent\textbf{Part III --- The audit.}
\begin{enumerate}[resume]
\item What did the firm report, and what does its usage log show in the first and last months?
\item What are the monthly fees used to price the shortfall?
\item What is owed in the first and last months?
\item What is owed over the three years, for display and for non-display?
\item Why does the audit cost the firm more than the back-billed fees?
\end{enumerate}

\medskip\noindent\textbf{Part IV --- The verdict.}
\begin{enumerate}[resume]
\item State the \emph{named result}: the licence shortfall in display users and non-display applications found by reconciling three years of
usage logs, and its back-billed cost under the venue's 2026 fees.
\item How long must usage logs be kept, and in what unit?
\item Is a conflated quote \omterm{def:pl:market-data-distribution-and-entitlements:derived}{derived data}?
\item What would you change first in the firm's platform?
\item In one sentence: what does a market-data licence pay for?
\end{enumerate}
\end{problem}

\section{Interview questions}

\begin{interviewq}[$\star$ \iqroles{developer}]\label{iq:pl:market-data-distribution-and-entitlements:1}
What is conflation, and when would you use it?
\end{interviewq}

\begin{interviewq}[$\star\star$ \iqroles{developer}]\label{iq:pl:market-data-distribution-and-entitlements:2}
A slow consumer is holding up your market-data bus. What do you do?
\end{interviewq}

\begin{interviewq}[$\star\star$ \iqroles{developer}]\label{iq:pl:market-data-distribution-and-entitlements:3}
What is the difference between display and non-display use of market data, and why does it matter to a developer?
\end{interviewq}

\begin{interviewq}[$\star\star$ \iqroles{developer}]\label{iq:pl:market-data-distribution-and-entitlements:4}
An exchange announces an audit of your data usage over the past three years. What do you need to have?
\end{interviewq}

\begin{interviewq}[$\star\star\star$ \iqroles{developer}]\label{iq:pl:market-data-distribution-and-entitlements:5}
Design the market-data distribution and \omterm{def:pl:market-data-distribution-and-entitlements:entitlement}{entitlement} platform of a firm with a thousand users and a hundred applications.
\end{interviewq}
